\documentclass[10pt,a4paper]{amsart}
\usepackage[margin=19mm]{geometry}
\usepackage{amsmath,amssymb,mathtools}
\usepackage[scr=rsfs,cal=euler]{mathalfa}
\usepackage{xfrac}
\usepackage[unicode,hidelinks,bookmarks=false]{hyperref}
\newcommand{\ie}{i.e.\ }
\newcommand{\cf}{cf.\ }
\newcommand{\resp}{respectively\ }
\newcommand{\Subsection}{\subsection}
\providecommand{\nobreakdash}{\nobreak\hbox{-}}
\setlength{\emergencystretch}{2em}
\multlinegap0pt
\usepackage{mathtools}
\usepackage{amssymb}
\newtheorem{theorem}{Theorem}
\newcommand{\Img}{{\rm Img}}
\newcommand{\Ker}{{\rm Ker}}
\newcommand{\ba}{\mathbf{a}}
\newcommand{\bff}{\mathbf{f}}
\newcommand{\bh}{\mathbf{h}}
 \newcommand{\bn}{\mathbf{n}}
\newcommand{\bv}{\mathbf{v}}
\newcommand{\curl}{\operatorname{curl}}
\newcommand{\divv}{\operatorname{div}}
\newcommand{\dx}{\operatorname{\mathbf{dx}}}
\title{On the Multi-Dimensional Divergence-Curl Problem and Its Connection with Pseudo-Harmonic Fields}
\author{A.V. Gorshkov}
\date{Source: arXiv:2605.10422v1 (CC BY 4.0)}
\begin{document}
\maketitle

\noindent\emph{Source and licence.} On the Multi-Dimensional Divergence-Curl Problem and Its Connection with Pseudo-Harmonic Fields. Version arXiv:2605.10422v1 (CC BY 4.0), distributed by arXiv under the Creative Commons Attribution 4.0 International licence, http://creativecommons.org/licenses/by/4.0/, as recorded in the arXiv licence metadata for this exact version and shown on the abstract page https://arxiv.org/abs/2605.10422v1. Full licence notice: \texttt{licenses/arxiv-2605.10422-CC-BY-4.0.txt}.\par\smallskip

\noindent\textit{Editorial scope.} Counted unit: the article's theorem mainTh (source0.tex lines 282--288). The article's complete original proof of it is delivered with it (source0.tex lines 290--310, one \texttt{proof} environment of 21 lines). No mathematics is written, repaired or re-derived: the statement and the proof are the article's own lines, in the article's own notation, and no retained line is substituted or re-flowed.  Retained uncounted as the source-local context the proof needs: lines 71--76; lines 78--80; lines 270--274.  Nothing was omitted from any retained span, so the package carries no omission record.  No retained line contains a citation, so the wrapper carries no bibliography and no bibliography entry.  No retained line contains a figure environment, an image-inclusion command or drawing code, so no image is delivered and the package declares no asset dependency.  Editorial formatting only, in addition: an AMS \texttt{amsart} preamble that restates the article's own declarations and macros used by the retained lines, namely the environments \texttt{theorem} on separate counters, together with the standard packages those lines need.  The retained environments print in this document as Theorem 1 in order of appearance, whereas the article prints the same items with the numbering of its own full document.  The complete native source of the article as distributed in the arXiv e-print for this exact version is bundled unchanged as \texttt{sources/200295/source0.tex}.
\begin{align}\label{mainsyst}
\begin{cases}
\curl \bv = \bff,\\
\divv \bv = 0,
\end{cases}
\end{align}

\begin{align}\label{boundcond}
\bv|_{\partial \Omega} = 0.
\end{align}

\begin{align}
&\curl^2 \ba = \bff, \label{maineq1} \\
&\curl \ba  \times \bn |_{\partial \Omega} = 0,\label{mainsystcond1}\\
&\ba  \times \bn |_{\partial \Omega} = 0.\label{mainsystcond2}
\end{align}

\begin{theorem}\label{mainTh}
Let $\Omega$ be a bounded Lipschitz domain and $\bff\in L_2(\Omega)$. The problem (\ref{mainsyst}), (\ref{boundcond}) has a unique solution if and only if:
\begin{enumerate}
\item $\divv \bff = 0,$
\item $\int_{\Omega} (\bff,\bh)\dx=0~\forall \bh \in \Ker \left ( \curl^2 \right ).$
\end{enumerate}
\end{theorem}

\begin{proof}
Solenoidality of $\bff$ follows from properties of the exterior derivative. Necessity of the orthogonality condition follows from the definition of the solution and the generalized Stokes formula. Multiplying (\ref{maineq1}) by $\bh$ with $\curl^2 \bh =0$ and integrating by parts, all boundary terms vanish due to the boundary conditions, giving
\begin{align*}
\int_M (\curl^2 \ba, \bh) \dx = 
\int_M (\ba, \curl^2 \bh) \dx +
\int_{\partial M} (\curl \ba \times \bn, \bh)\,ds - 
\int_{\partial M} (\ba \times \bn, \curl \bh) = 0
\end{align*}
for all $\bh \in \Ker \left ( \curl^2 \right )$.

For sufficiency: the operator $\curl^2: H_0(\curl^2,\Omega)\to L_2(\Omega)$ has an adjoint $\curl^{2,*}: L_2(\Omega) \to H^*(\Omega)$. By the boundary conditions, $\curl^2$ and $\curl^{2,*}$ agree on smooth functions, and $\Img \curl^2 = \left ( \Ker \curl^{2,*} \right )^\perp$. Since $\curl^{2,*}$ is an epimorphism, for any $\bff \in L_2$ orthogonal to $\Ker \curl^{2,*}$, a solution $\ba \in H_0(\curl^2,\Omega)$ exists.

Uniqueness follows from Stokes' formula. Let $\ba$ be a solution with $\bff=0$. Testing against $\ba$:
\begin{align*}
\int_\Omega (\curl^2 \ba,\ba) \dx= 
\int_\Omega (\curl \ba, \curl \ba) \dx + 
\int_{\partial \Omega} (\curl \ba \times \bn, \ba) \dx =
\int_\Omega (\curl \ba, \curl \ba) \dx = \\ \|\curl \ba\|_{L_2} = 0.
\end{align*}
Hence $\bv = \curl \ba \equiv 0$.  
\end{proof}

\end{document}
